% chapters/01_einleitung.tex

\chapter{Einleitung}
\label{chap:einleitung}

\section{Motivation}
\label{sec:motivation}

Humanoide Roboter stehen vor der Herausforderung, komplexe Manipulations-
aufgaben in unstrukturierten Umgebungen auszuführen.
Traditionelle Programmieransätze erfordern aufwändige manuelle Regelentwicklung
und sind nur begrenzt auf neue Aufgaben übertragbar.
Foundation-Modelle aus dem Bereich der künstlichen Intelligenz bieten einen
vielversprechenden Ansatz, diesen Prozess durch \ac{il} zu vereinfachen:
Ein einmal vortrainiertes \ac{vla}-Modell kann durch gezieltes Fine-Tuning auf
wenigen Demonstrations-Episoden an neue Roboter und Aufgaben angepasst werden.

\begin{infobox}[Projektziel]
  Diese Projektarbeit untersucht, ob und wie das Modell NVIDIA GR00T N1.6 auf dem
  Unitree G1 Humanoidroboter mit DEX3-Hand für Block-Stacking-Aufgaben mittels
  Imitation Learning erfolgreich angepasst werden kann.
\end{infobox}

\section{Problemstellung}
\label{sec:problemstellung}

Das Block-Stacking mit einer humanoiden Hand stellt eine nicht-triviale
Manipulation dar: Der Roboter muss Objekte präzise greifen, heben und stapeln,
wobei Kamerasicht, Propriozeption und Kraft-Feedback in einer gemeinsamen
Repräsentation verarbeitet werden müssen.
Folgende Herausforderungen ergeben sich:

\begin{itemize}
  \item Hohe \ac{dof}-Anzahl des Systems (G1-Körper + DEX3-Hand)
  \item Begrenzte Menge an Demonstrations-Daten
  \item Rechenintensives Training auf \ac{gpu}-Clustern
  \item Sim-to-Real-Gap bei der Evaluation
\end{itemize}

\section{Zielsetzung}
\label{sec:zielsetzung}

Ziele dieser Arbeit sind:

\begin{enumerate}
  \item Aufbau einer reproduzierbaren Trainingsinfrastruktur (Docker / Apptainer)
  \item Fine-Tuning von GR00T N1.6 auf dem G1-Datensatz mittels LeRobot-Format
  \item Quantitative Evaluation des trainierten Modells
  \item Dokumentation und Übertragbarkeit der Pipeline auf weitere Aufgaben
\end{enumerate}

\section{Aufbau der Arbeit}
\label{sec:aufbau}

\Cref{chap:grundlagen} legt die theoretischen Grundlagen zu humanoider Robotik,
\acp{vla} und \ac{il} dar.
\Cref{chap:methoden} beschreibt den Datensatz, die Systemarchitektur und den
Trainingsprozess.
In \cref{chap:ergebnisse} werden die Trainings- und Evaluationsergebnisse
vorgestellt.
\Cref{chap:diskussion} diskutiert die Ergebnisse und Limitierungen.
\Cref{chap:fazit} schließt die Arbeit mit einem Fazit und Ausblick ab.
